\documentclass[12pt, a4paper]{article}
\usepackage{liturg}
\usepackage[latin,english]{babel}
\usepackage{ecclesiastic}
\usepackage[utf8]{inputenc}
\title{Missa Præsanctificatorum}
\usepackage{fontspec}
\usepackage{blindtext}
\date{}


\begin{document}
	\begin{titlepage}
		\centering
		\vspace*{3cm}
		{\Huge\bfseries MISSA PRÆSANCTIFICATORUM\par}
		\vspace{2cm}
		{\Large Mass of the Presanctified\par}
		\vfill
	\end{titlepage}
	
	\section{Lectiones}
	
	\instruct{In Choro, dicta Nona, Sacerdos et Ministri induti paramentis nigri coloris, sine luminaribus et incenso, procedunt ad Altare: et ante illud prostrati aliquamdiu orant. Interim Acolythi unam tantum tobaleam extendunt super Altare. Sacerdos cum Ministris facta oratione ascendit ad Altare, et osculatur illud in medio: deinde Lector accedit ad legendum Prophetiam in loco, ubi legitur Epistola, et incipit eam sine titulo: quam etiam Sacerdos legit submissa voce apud Altare in cornu Epistolæ}
	\\
	\instruct{Osee 6, 1-6.}
	Hæc dicit Dóminus: In tribulatióne sua mane consúrgent ad me: Veníte, et revertámur ad Dóminum: quia ipse cepit, et sanábit nos: percútiet, et curábit nos. Vivificábit nos post duos dies: in die tértia suscitábit nos, et vivémus in conspéctu ejus. Sciémus, sequemúrque, ut cognoscámus Dóminum: quasi dilúculum præparátus est egréssus ejus, et véniet quasi imber nobis temporáneus et serótinus terræ.
	\\\\
	Quid fáciam tibi, Ephraim? Quid fáciam tibi, Juda? misericórdia vestra quasi nubes matutína: et quasi ros mane pertránsiens. Propter hoc dolávi in prophétis, occídi eos in verbis oris mei: et judícia tua quasi lux egrediéntur. Quia misericórdiam vólui, et non sacrifícium, et sciéntiam Dei, plus quam holocáusta.
	\\\\
	\instruct{Et non respondetur Deo gratias, quod servatur etiam post Lectionem sequentem.}
	\\
	\instruct{Tractus -- Habacuc 3.}
	Dómine, audívi audítum tuum, et tímui: considerávi ópera tua, et expávi.
	\\
	\priest{In médio duórum animálium innotescéris: dum appropinquáverint anni, sognoscéris: dum advénerit tempus, osténdens,}
	\priest{In eo, dum conturbáta fúerit ánima mea: in ira, misericórdiæ memor eris.}
	\priest{Deus a Líbano véniet, et Sanctus de monte umbróso et condénso.}
	\priest{Opéruit cœlos majéstas ejus: et laudis ejus plena est terra.}
	\\
	\instruct{Finito Tractu, Sacerdos in cornu Epistolæ dicit:}
	\\
	Orémus
	\\
	\priest{Flectámus génua.}
	\server{Leváte.}
	\\
	Deus, a quo et Judas reátus sui pœnam, et confessiónis suæ latro præmium sumpsit, concéde nobis tuæ propitiatiónis efféctum: ut, sicut in passióne sua Jesus Christus, Dóminus noster, divérsa utrísque íntulit stipéndia meritórum; ita nobis, abláto vetustátis erróre, resurrectiónis suæ grátiam largiátur:
	\\\\
	Qui tecum vivit et regnat in unitáte Spíritus Sancti Deus per ómnia sæcula sæculórum.
	\\\\
	\server{Amen.}
	\\
	\instruct{Subdiaconus in tono Epistolæ similiter sine titulo cantat sequentem Lectionem:}
	\\
	\instruct{Exodi 12, 1-11.}
	In diébus illis: Dixit Dóminus ad Móysen et Aaron in terra Ægýpti: Mensis iste vobis princípium ménsium primus erit in ménsibus anni. Loquímini ad univérsum cœtum filiórum Israël, et dícite eis: Décima die mensis hujus tollat unusquísque agnum per famílias et domos suas. Sin autem minor est númerus, ut suffícere possit ad vescéndum agnum, assúmet vicínum suum, qui junctus est dómui suæ, juxta númerum animárum, quæ suffícere possunt ad esum agni. Erit autem agnus absque mácula, másculus, annículus: juxta quem ritum tollétis et hædum.
	\\\\
	Et servábitis eum usque ad quartam décimam diem mensis hujus: immolabítque eum univérsa multitúdo filiórum Israël ad vésperam. Et sument de sánguine ejus, ac ponent super utrúmque postem et in superlimináribus domórum, in quibus cómedent illum. Et edent carnes nocte illa assas igni, et ázymos panes cum lactúcis agréstibus. Non comedétis ex eo crudum quid nec coctum aqua, sed tantum assum igni: caput cum pédibus ejus et intestínis vorábitis. Nec remanébit quidquam ex eo usque mane. Si quid resíduum fúerit, igne comburétis.
	\\\\
	Sic autem comedétis illum: Renes vestros accingétis, et calceaménta habébitis in pédibus, tenéntes báculos in mánibus, et comedétis festinánter: est enim Phase id est tránsitus Dómini.
	\\\\
	\instruct{Tractus -- Ps 139:2-10 et 14.}
	Eripe me, Dómine, ab hómine malo: a viro iníquo líbera me.
	\\
	\priest{Qui cogitavérunt malítias in corde: tota die constituébant prælia.}
	\priest{Acuérunt linguas suas sicut serpéntis: venénum áspidum sub lábiis eórum. Custódi me, Dómine, de manu peccatóris: et ab homínibus iníquis líbera me.}
	\priest{Qui cogitavérunt supplantáre gressus meos: abscondérunt supérbi láqueum mihi.}
	\priest{Et funes extendérunt in láqueum pédibus meis: juxta iter scándalum posuérunt mihi.}
	\priest{Dixi Dómino: Deus meus es tu: exáudi, Dómine, vocem oratiónis meæ.}
	\priest{Dómine, Dómine, virtus salútis meæ: obúmbra caput meum in die belli.}
	\priest{Ne tradas me a desidério meo peccatóri: cogitavérunt advérsus me: ne derelínquas me, ne umquam exalténtur.}
	\priest{Caput circúitus eórum: labor labiórum ipsórum opériet eos.}
	\priest{Verúmtamen justi confitebúntur nómini tuo: et habitábunt recti cum vultu tuo.}
	
	\section{Passio}
	
	\instruct{Finito Tractu, dicitur Passio super nudum pulpitum: quam Celebrans submissa voce legit in cornu Epistolæ.}
	\\
	Pássio Dómini nostri Jesu Christi secúndum Joánnem
	\\\\
	\instruct{Joann. 18:1-40; 19:1-42.}
	In illo témpore: Egréssus est Jesus cum discípulis suis trans torréntem Cedron, ubi erat hortus, in quem introívit ipse et discípuli ejus. Sciébat autem et Judas, qui tradébat eum, locum: quia frequénter Jesus convénerat illuc cum discípulis suis. Judas ergo cum accepísset cohórtem, et a pontifícibus et pharisæis minístros, venit illuc cum latérnis et fácibus et armis. Jesus ítaque sciens ómnia, quæ ventúra erant super eum, procéssit, et dixit eis:
	\\\\
	\priest{Quem quæritis?}
	\\
	Respondérunt ei:
	\\\\
	\server{Jesum Nazarénum.}
	\\
	Dicit eis Jesus:
	\\\\
	\priest{Ego sum.}
	\\
	Stabat autem et Judas, qui tradébat eum, cum ipsis. Ut ergo dixit eis: Ego sum: abiérunt retrórsum, et cecidérunt in terram. Iterum ergo interrogávit eos:
	\\\\
	\priest{Quem quæritis?}
	\\
	Illi autem dixérunt:
	\\\\
	\server{Jesum Nazarénum.}
	\\
	Respóndit Jesus:
	\\\\
	\priest{Dixi vobis, quia ego sum: si ergo me quæritis, sínite hos abíre.}
	\\
	Ut implerétur sermo, quem dixit: Quia quos dedísti mihi, non pérdidi ex eis quemquam. Simon ergo Petrus habens gládium edúxit eum: et percússit pontíficis servum: et abscídit aurículam ejus déxteram. Erat autem nomen servo Malchus. Dixit ergo Jesus Petro:
	\\\\
	\priest{Mitte gládium tuum in vagínam. Cálicem, quem dedit mihi Pater, non bibam illum?}
	\\
	Cohors ergo et tribúnus et minístri Judæórum comprehendérunt Jesum, et ligavérunt eum: et adduxérunt eum ad Annam primum, erat enim socer Cáiphæ, qui erat póntifex anni illíus. Erat autem Cáiphas, qui consílium déderat Judæis: Quia éxpedit, unum hóminem mori pro pópulo. Sequebátur autem Jesum Simon Petrus et álius discípulus. Discípulus autem ille erat notus pontífici, et introívit cum Jesu in átrium pontíficis. Petrus autem stabat ad óstium foris. Exívit ergo discípulus álius, qui erat notus pontífici, et dixit ostiáriæ: et introdúxit Petrum. Dicit ergo Petro ancílla ostiária: 
	\\\\
	\server{Numquid et tu ex discípulis es hóminis istíus?}
	\\
	Dicit ille:
	\\\\
	\server{Non sum.}
	\\
	Stabant autem servi et minístri ad prunas, quia frigus erat, et calefaciébant se: erat autem cum eis et Petrus stans et calefáciens se. Póntifex ergo interrogávit Jesum de discípulis suis et de doctrína ejus. Respóndit ei Jesus: 
	\\\\
	\priest{Ego palam locútus sum mundo: ego semper dócui in synagóga et in templo, quo omnes Judæi convéniunt: et in occúlto locútus sum nihil. Quid me intérrogas? Intérroga eos, qui audiérunt, quid locútus sim ipsis: ecce, hi sciunt, quæ díxerim ego.}
	\\
	Hæc autem cum dixísset, unus assístens ministrórum dedit álapam Jesu, dicens:
	\\\\
	\server{Sic respóndes pontífici?}
	\\
	Respóndit ei Jesus:
	\\\\
	\priest{Si male locútus sum, testimónium pérhibe de malo: si autem bene, quid me cædis?}
	\\
	Et misit eum Annas ligátum ad Cáipham pontíficem. Erat autem Simon Petrus stans et calefáciens se. Dixérunt ergo ei:
	\\\\
	\server{Numquid et tu ex discípulis ejus es?}
	\\
	Negávit ille et dixit:
	\\\\
	\server{Non sum.}
	\\
	Dicit ei unus ex servis pontíficis, cognátus ejus, cujus abscídit Petrus aurículam:
	\\\\
	\server{Nonne ego te vidi in horto cum illo?}
	\\
	Iterum ergo negávit Petrus: et statim gallus cantávit. Addúcunt ergo Jesum a Cáipha in prætórium. Erat autem mane: et ipsi non introiérunt in prætórium, ut non contaminaréntur, sed ut manducárent pascha. Exívit ergo Pilátus ad eos foras et dixit: 
	\\\\
	\server{Quam accusatiónem affértis advérsus hóminem hunc? }
	\\
	Respondérunt et dixérunt ei:
	\\\\
	\server{Si non esset hic malefáctor, non tibi tradidissémus eum.}
	\\
	Dixit ergo eis Pilátus: 
	\\\\
	\server{Accípite eum vos, et secúndum legem vestram judicáte eum.}
	\\
	Dixérunt ergo ei Judæi:
	\\\\
	\server{Nobis non licet interfícere quemquam.}
	\\
	Ut sermo Jesu implerétur, quem dixit, signíficans, qua morte esset moritúrus. Introívit ergo íterum in prætórium Pilátus, et vocávit Jesum et dixit ei:
	\\\\
	\server{Tu es Rex Judæórum?}
	\\
	Respóndit Jesus:
	\\\\
	\priest{A temetípso hoc dicis, an álii dixérunt tibi de me?}
	\\
	Respóndit Pilátus:
	\\\\
	\server{Numquid ego Judæus sum? Gens tua et pontífices tradidérunt te mihi: quid fecísti?}
	\\
	Respóndit Jesus:
	\\\\
	\priest{Regnum meum non est de hoc mundo. Si ex hoc mundo esset regnum meum, minístri mei útique decertárent, ut non tráderer Judæis: nunc autem regnum meum non est hinc.}
	\\
	Dixit ítaque ei Pilátus:
	\\\\
	\server{Ergo Rex es tu?}
	\\
	Respóndit Jesus:
	\\\\
	\priest{Tu dicis, quia Rex sum ego. Ego in hoc natus sum et ad hoc veni in mundum, ut testimónium perhíbeam veritáti: omnis, qui est ex veritáte, audit vocem meam.}
	\\
	Dicit ei Pilátus:
	\\\\
	\server{Quid est véritas?}
	\\
	Et cum hoc dixísset, íterum exívit ad Judæos, et dicit eis:
	\\\\
	\server{Ego nullam invénio in eo causam. Est autem consuetúdo vobis, ut unum dimíttam vobis in Pascha: vultis ergo dimíttam vobis Regem Judæórum?}
	\\
	Clamavérunt ergo rursum omnes, dicéntes:
	\\\\
	\server{Non hunc, sed Barábbam.}
	\\
	Erat autem Barábbas latro. Tunc ergo apprehéndit Pilátus Jesum et flagellávit. Et mílites plecténtes corónam de spinis, imposuérunt cápiti ejus: et veste purpúrea circumdedérunt eum. Et veniébant ad eum, et dicébant:
	\\\\
	\server{Ave, Rex Judæórum.}
	\\
	Et dabant ei álapas. Exívit ergo íterum Pilátus foras et dicit eis:
	\\\\
	\server{Ecce, addúco vobis eum foras, ut cognoscátis, quia nullam invénio in eo causam.}
	\\
	Exívit ergo Jesus portans corónam spíneam et purpúreum vestiméntum. Et dicit eis:
	\\\\
	\server{Ecce homo.}
	\\
	Cum ergo vidíssent eum pontífices et minístri, clamábant, dicéntes:
	\\\\
	\server{Crucifíge, crucifíge eum.}
	\\
	Dicit eis Pilátus:
	\\\\
	\server{Accípite eum vos et crucifígite: ego enim non invénio in eo causam.}
	\\
	Respondérunt ei Judæi:
	\\\\
	\server{Nos legem habémus, et secúndum legem debet mori, quia Fílium Dei se fecit.}
	\\
	Cum ergo audísset Pilátus hunc sermónem, magis tímuit. Et ingréssus est prætórium íterum: et dixit ad Jesum:
	\\\\
	\server{Unde es tu?}
	\\
	Jesus autem respónsum non dedit ei. Dicit ergo ei Pilátus:
	\\\\
	\server{Mihi non loquéris? Nescis, quia potestátem hábeo crucifígere te, et potestátem hábeo dimíttere te?}
	\\
	Respóndit Jesus:
	\\\\
	\priest{Non habéres potestátem advérsum me ullam, nisi tibi datum esset désuper. Proptérea, qui me trádidit tibi, majus peccátum habet.}
	\\
	Et exínde quærébat Pilátus dimíttere eum. Judæi autem clamábant dicéntes: 
	\\\\
	\server{Si hunc dimíttis, non es amícus Cæsaris. Omnis enim, qui se regem facit, contradícit Cæsari.}
	\\
	Pilátus autem cum audísset hos sermónes, addúxit foras Jesum, et sedit pro tribunáli, in loco, qui dícitur Lithóstrotos, hebráice autem Gábbatha. Erat autem Parascéve Paschæ, hora quasi sexta, et dicit Judæis:
	\\\\
	\server{Ecce Rex vester.}
	\\
	Illi autem clamábant:
	\\\\
	\server{Tolle, tolle, crucifíge eum.}
	\\
	Dicit eis Pilátus:
	\\\\
	\server{Regem vestrum crucifígam?}
	\\
	Respondérunt pontífices:
	\\\\
	\server{Non habémus regem nisi Cæsarem.}
	\\
	Tunc ergo trádidit eis illum, ut crucifigerétur. Suscepérunt autem Jesum et eduxérunt. Et bájulans sibi crucem, exívit in eum, qui dícitur Calváriæ, locum, hebráice autem Gólgotha: ubi crucifixérunt eum, et cum eo álios duos, hinc et hinc, médium autem Jesum. Scripsit autem et títulum Pilátus: et pósuit super crucem. Erat autem scriptum: Jesus Nazarénus, Rex Judæórum. Hunc ergo títulum multi Judæórum legérunt, quia prope civitátem erat locus, ubi crucifíxus est Jesus. Et erat scriptum hebráice, græce et latíne. Dicébant ergo Piláto pontífices Judæórum:
	\\\\
	\server{Noli scríbere Rex Judæórum, sed quia ipse dixit: Rex sum Judæórum.}
	\\
	Respóndit Pilátus:
	\\\\
	\server{Quod scripsi, scripsi.}
	\\
	Mílites ergo cum crucifixíssent eum, accepérunt vestiménta ejus et fecérunt quátuor partes: unicuíque míliti partem, et túnicam. Erat autem túnica inconsútilis, désuper contéxta per totum. Dixérunt ergo ad ínvicem:
	\\\\
	\server{Non scindámus eam, sed sortiámur de illa, cujus sit.}
	\\
	Ut Scriptúra implerétur, dicens: Partíti sunt vestiménta mea sibi: et in vestem meam misérunt sortem. Et mílites quidem hæc fecérunt. Stabant autem juxta crucem Jesu Mater ejus et soror Matris ejus, María Cléophæ, et María Magdaléne. Cum vidísset ergo Jesus Matrem et discípulum stantem, quem diligébat, dicit Matri suæ:
	\\\\
	\priest{Múlier, ecce fílius tuus.}
	\\
	Deínde dicit discípulo:
	\\\\
	\priest{Ecce mater tua.}
	\\
	Et ex illa hora accépit eam discípulus in sua. Póstea sciens Jesus, quia ómnia consummáta sunt, ut consummarétur Scriptúra, dixit:
	\\\\
	\priest{Sítio.}
	\\
	Vas ergo erat pósitum acéto plenum. Illi autem spóngiam plenam acéto, hyssópo circumponéntes, obtulérunt ori ejus. Cum ergo accepísset Jesus acétum, dixit:
	\\\\
	\priest{Consummátum est.}
	\\
	Et inclináto cápite trádidit spíritum.
	\\\\
	\instruct{Hic genuflectitur, et pausatur aliquantulum}
	\\
	Judæi ergo - quóniam Parascéve erat - ut non remanérent in cruce córpora sábbato - erat enim magnus dies ille sábbati - rogavérunt Pilátum, ut frangeréntur eórum crura et tolleréntur. Venérunt ergo mílites: et primi quidem fregérunt crura et altérius, qui crucifíxus est cum eo. Ad Jesum autem cum veníssent, ut vidérunt eum jam mórtuum, non fregérunt ejus crura, sed unus mílitum láncea latus ejus apéruit, et contínuo exívit sanguis et aqua. Et qui vidit, testimónium perhíbuit: et verum est testimónium ejus. Et ille scit, quia vera dicit: ut et vos credátis. Facta sunt enim hæc, ut Scriptúra implerétur: Os non comminuétis ex eo. Et íterum ália Scriptúra dicit: Vidébunt in quem transfixérunt.
	\\\\
	\instruct{Quod sequitur, legitur in tono Evangelii: et dicitur Munda cor meum, sed non petitur benedíctio, et non deferuntur luminaria neque incensum, et Celebrans in fine non osculatur librum.}
	\\
	Munda cor meum ac lábia mea, omnípotens Deus, qui lábia Isaíæ Prophétæ cálculo mundásti igníto: ita me tua grata miseratióne dignáre mundáre, ut sanctum Evangélium tuum digne váleam nuntiáre. Per Christum, Dóminum nostrum. Amen.
	\\\\
	\priest{Dóminus vobíscum.}
	\server{Et cum spíritu tuo.}
	\\
	Post hæc autem rogávit Pilátum Joseph ab Arimathæa - eo quod esset discípulus Jesu - occúltus autem propter metum Judæórum, ut tólleret corpus Jesu. Et permísit Pilátus. Venit ergo et tulit corpus Jesu. Venit autem et Nicodémus, qui vénerat ad Jesum nocte primum, ferens mixtúram myrrhæ et áloës, quasi libras centum. Accepérunt ergo corpus Jesu, et ligavérunt illud línteis cum aromátibus, sicut mos est Judæis sepelíre. Erat autem in loco, ubi crucifíxus est, hortus: et in horto monuméntum novum, in quo nondum quisquam pósitus erat. Ibi ergo propter Parascéven Judæórum, quia juxta erat monuméntum, posuérunt Jesum.
	
	\section{Intercessiones}
	
	\instruct{Deinde Sacerdos stans in cornu Epistolæ, incipit absolute junctis manibus:}
	\\
	\instruct{Pro Ecclesia}
	Orémus, dilectíssimi nobis, pro Ecclésia sancta Dei: ut eam Deus et Dóminus noster pacificáre, adunáre, et custodíre dignétur toto orbe terrárum: subjíciens ei principátus et potestátes: detque nobis quietam et tranquíllam vitam degentibus, glorificáre Deum, Patrem omnipoténtem.
	\\\\
	Orémus
	\\
	\server{Flectámus génua.}
	\server{Leváte.}
	\\
	Omnípotens sempitérne Deus, qui glóriam tuam ómnibus in Christo géntibus revelásti: custódi ópera misericórdiæ tuæ; ut Ecclésia tua, toto orbe diffúsa, stábili fide in confessióne tui nóminis persevéret.
	\\\\
	Per eúndem Dóminum nostrum Jesum Christum Fílium tuum, qui tecum vivit et regnat in unitáte Spíritus Sancti, Deus, per ómnia sæcula sæculórum.
	\\\\
	\server{Amen.}
	\\
	\instruct{Pro Papa}
	Orémus et pro beatíssimo Papa nostro N... ut Deus et Dóminus noster, qui elégit eum in órdine episcopátus, salvum atque incólumem custódiat Ecclésiæ suæ sanctæ, ad regéndum pópulum sanctum Dei.
	\\\\
	Orémus
	\\
	\server{Flectámus génua.}
	\server{Leváte.}
	\\
	Deus, cujus judício univérsa fundántur: réspice propítius ad preces nostras, et electum nobis Antístitem tua pietáte consérva; ut christiána plebs, quæ te gubernátur auctóre, sub tanto Pontífice, credulitátis suæ méritis augeátur.
	\\\\
	Per Dóminum nostrum Jesum Christum Fílium tuum, qui tecum vivit et regnat in unitáte Spíritus Sancti, Deus, per ómnia sæcula sæculórum.
	\\\\
	\server{Amen.}
	\\
	\instruct{Pro clerici et populo}
	Orémus et pro ómnibus Epíscopis, Presbýteris, Diacónibus, Subdiacónibus, Acólythis, Exorcístis, Lectóribus, Ostiáriis, Confessóribus, Virgínibus, Víduis: et pro omni pópulo sancto Dei.
	\\\\
	Orémus
	\\
	\server{Flectámus génua.}
	\server{Leváte.}
	\\
	Deus, cujus Spíritu totum corpus Ecclésiæ sanctificátur et régitur: exáudi nos pro univérsis ordínibus supplicántes; ut, grátiæ tuæ múnere, ab ómnibus tibi grádibus fidéliter serviátur.
	\\\\
	Per Dóminum nostrum Jesum Christum Fílium tuum, qui tecum vivit et regnat in unitáte Spíritus Sancti, Deus, per ómnia sæcula sæculórum.
	\\\\
	\server{Amen.}
	\\
	\instruct{Pro catechumenis}
	Orémus et pro catechúmenis nostris: ut Deus et Dóminus noster adapériat aures præcordiórum ipsórum januámque misericordiæ; ut, per lavácrum regeneratiónis accépta remissióne ómnium peccatórum, et ipsi inveniántur in Christo Jesu, Dómino nostro.
	\\\\
	Orémus
	\\
	\server{Flectámus génua.}
	\server{Leváte.}
	\\
	Omnípotens sempitérne Deus, qui Ecclésiam tuam nova semper prole fecúndas: auge fidem et intellectum catechúmenis nostris; ut, renáti fonte baptismátis, adoptiónis tuæ fíliis aggregéntur.
	\\\\
	Per Dóminum nostrum Jesum Christum Fílium tuum, qui tecum vivit et regnat in unitáte Spíritus Sancti, Deus, per ómnia sæcula sæculórum.
	\\\\
	\server{Amen.}
	\\
	\instruct{Pro contra errores}
	Orémus, dilectíssimi nobis, Deum Patrem omnipoténtem, ut cunctis mundum purget erróribus: morbos áuferat: famem depéllat: apériat cárceres: víncula dissólvat: peregrinántibus réditum: infirmántibus sanitátem: navigántibus portum salútis indúlgeat.
	\\\\
	Orémus
	\\
	\server{Flectámus génua.}
	\server{Leváte.}
	\\
	Omnípotens sempitérne Deus, mæstórum consolátio, laborántium fortitúdo: pervéniant ad te preces de quacúmque tribulatióne clamántium; ut omnes sibi in necessitátibus suis misericórdiam tuam gáudeant affuísse.
	\\\\
	Per Dóminum nostrum Jesum Christum Fílium tuum, qui tecum vivit et regnat in unitáte Spíritus Sancti, Deus, per ómnia sæcula sæculórum.
	\\\\
	\server{Amen.}
	\\
	\instruct{Pro hæreticis et schismaticis}
	Orémus et pro hæréticis et schismáticis: ut Deus et Dóminus noster éruat eos ab erróribus univérsis; et ad sanctam matrem Ecclésiam Cathólicam atque Apostólicam revocáre dignétur.
	\\\\
	Orémus
	\\
	\server{Flectámus génua.}
	\server{Leváte.}
	\\
	Omnípotens sempitérne Deus, qui salvas omnes, et néminem vis períre: réspice ad ánimas diabólica fraude decéptas; ut, omni hærética pravitáte depósita, errántium corda resipíscant, et ad veritátis tuæ rédeant unitátem.
	\\\\
	Per Dóminum nostrum Jesum Christum Fílium tuum, qui tecum vivit et regnat in unitáte Spíritus Sancti, Deus, per ómnia sæcula sæculórum.
	\\\\
	\server{Amen.}
	\\
	\instruct{Pro Judæis}
	Orémus et pro pérfidis Judæis: ut Deus et Dóminus noster áuferat velámen de córdibus eórum; ut et ipsi agnóscant Jesum Christum, Dóminum nostrum.
	\\\\
	Omnípotens sempitérne Deus, qui étiam judáicam perfídiam a tua misericórdia non repéllis: exáudi preces nostras, quas pro illíus pópuli obcæcatióne deférimus; ut, ágnita veritátis tuæ luce, quæ Christus est, a suis ténebris eruántur.
	\\\\
	Per Dóminum nostrum Jesum Christum Fílium tuum, qui tecum vivit et regnat in unitáte Spíritus Sancti, Deus, per ómnia sæcula sæculórum.
	\\\\
	\server{Amen.}
	\\
	\instruct{Pro paganis}
	Orémus et pro pagánis: ut Deus omnípotens áuferat iniquitátem a córdibus eórum; ut, relíctis idólis suis, convertántur ad Deum vivum et verum, et únicum Fílium ejus Jesum Christum, Deum et Dóminum nostrum.
	\\\\
	Orémus
	\\
	\server{Flectámus génua.}
	\server{Leváte.}
	\\
	Omnípotens sempitérne Deus, qui non mortem peccatórum, sed vitam semper inquíris: súscipe propítius oratiónem nostram, et líbera eos ab idolórum cultúra; et ággrega Ecclésiæ tuæ sanctæ, ad laudem et glóriam nóminis tui.
	\\\\
	Per Dóminum nostrum Jesum Christum Fílium tuum, qui tecum vivit et regnat in unitáte Spíritus Sancti, Deus, per ómnia sæcula sæculórum.
	\\\\
	\server{Amen.}
	
	\section{Adoratio crucis}
	\instruct{Completis Orationibus, Sacerdos deposita Casula accedit ad cornu Epistolæ, et ibi in posteriori parte anguli Altaris accipit a Diacono Crucem jam in Altari præparatam: quam, versa facie ad populum, a summitate parum discooperit, incipiens solus Antiphonam Ecce lignum Crucis, ac deinceps in reliquis juvatur in cantu a Ministris usque ad Veníte, adorémus. Choro autem cantante Veníte, adorémus, omnes se prosternunt, excepto Celebrante.
	\\\\
	Deinde procedit ad anteriorem partem anguli ejusdem cornu Epistolæ: et discooperiens brachium dextrum Crucis, elevansque eam paulisper, altius quam primo, incipit: Ecce lignum Crucis, aliis cantantibus et adorantibus, ut supra. Deinde Sacerdos procedit ad medium Altaris: et discooperiens Crucem totaliter, ac elevans eam, tertio altius incipit: Ecce lignum Crucis, aliis cantantibus et adorantibus, ut supra.}
	\\
	\priest{Ecce lignum Crucis, in quo salus mundi pepéndit.}
	\server{Veníte, adorémus.}
	\\
	\priest{Ecce lignum Crucis, in quo salus mundi pepéndit.}
	\server{Veníte, adorémus.}
	\\
	\priest{Ecce lignum Crucis, in quo salus mundi pepéndit.}
	\server{Veníte, adorémus.}
	\\
	\instruct{Postea Sacerdos solus portat Crucem ad locum ante Altare præparatum, et genuflexus ibidem eam locat: mox depositis calceamentis, accedit ad adorandam Crucem, ter genua flectens antequam eam deosculetur. Hoc facto revertitur, et accipit calceamenta, et Casulam. Postmodum Ministri Altaris, deinde alii Clerici, et laici, bini et bini, ter genibus flexis, ut dictum est, Crucem adórant. Interim, dum fit adoratio Crucis, cantantur Improperia, et alia quæ sequuntur, vel omnia vel pars eorum, prout multitudo adorantium vel paucitas requirit: quæ etiam Sacerdos sedens ad scamnum legit cum Ministris, ut infra:}
	\\
	\priest{Pópule meus, quid feci tibi? aut in quo contristávi te? respónde mihi. Quia edúxi te de terra Ægýpti: parásti Crucem Salvatóri tuo.}
	\server{Agios o Theós. Sanctus Deus. Agios ischyrós. Sanctus fortis. Agios athánatos, eléison imas. Sanctus immortális, miserére nobis.}
	\\
	\priest{Quia edúxi te per desértum quadragínta annis, et manna cibávi te, et introdúxi te in terram satis bonam: parásti Crucem Salvatóri tuo.}
	\server{Agios o Theós. Sanctus Deus. Agios ischyrós. Sanctus fortis. Agios athánatos, eléison imas. Sanctus immortális, miserére nobis.}
	\\
	\priest{Quid ultra débui fácere tibi, et non feci? Ego quidem plantávi te víneam meam speciosíssimam: et tu facta es mihi nimis amára: acéto namque sitim meam potásti: et láncea perforásti latus Salvatóri tuo.}
	\server{Agios o Theós. Sanctus Deus. Agios ischyrós. Sanctus fortis. Agios athánatos, eléison imas. Sanctus immortális, miserére nobis.}
	\\
	\priest{Ego propter te flagellávi Ægýptum cum primogénitis suis: et tu me flagellátum tradidísti.}
	\server{Pópule meus, quid feci tibi? aut in quo contristávi te? respónde mihi.}
	\\
	\priest{Ego edúxi te de Ægýpto, demérso Pharaóne in Mare Rubrum: et tu me tradidísti princípibus sacerdótum.}
	\server{Pópule meus, quid feci tibi? aut in quo contristávi te? respónde mihi.}
	\\
	\priest{Ego ante te apérui mare: et tu aperuísti láncea latus meum.}
	\server{Pópule meus, quid feci tibi? aut in quo contristávi te? respónde mihi.}
	\\
	\priest{Ego ante te præívi in colúmna nubis: et tu me duxísti ad prætórium Piláti.}
	\server{Pópule meus, quid feci tibi? aut in quo contristávi te? respónde mihi.}
	\\
	\priest{Ego te pavi manna per desértum: et tu me cecidísti álapis et flagéllis.}
	\server{Pópule meus, quid feci tibi? aut in quo contristávi te? respónde mihi.}
	\\
	\priest{Ego te potávi aqua salútis de petra: et tu me potásti felle et acéto.}
	\server{Pópule meus, quid feci tibi? aut in quo contristávi te? respónde mihi.}
	\\
	\priest{Ego propter te Chananæórum reges percússi: et tu percussísti arúndine caput meum.}
	\server{Pópule meus, quid feci tibi? aut in quo contristávi te? respónde mihi.}
	\\
	\priest{Ego dedi tibi sceptrum regale: et tu dedísti capiti meo spíneam coronam.}
	\server{Pópule meus, quid feci tibi? aut in quo contristávi te? respónde mihi.}
	\\
	\priest{Ego te exaltávi magna virtúte: et tu me suspendísti in patíbulo Crucis.}
	\server{Pópule meus, quid feci tibi? aut in quo contristávi te? respónde mihi.}
	\\
	\priest{Crucem tuam adorámus, Dómine: et sanctam resurrectiónem tuam laudámus et glorificámus: ecce enim, propter lignum venit gaudium in univérso mundo.}
	\\
	\instruct{Ps 66:2.}
	\priest{Deus misereátur nostri et benedícat nobis: Illúminet vultum suum super nos et misereátur nostri.}
	\server{Crucem tuam adorámus, Dómine: et sanctam resurrectiónem tuam laudámus et glorificámus: ecce enim, propter lignum venit gáudium in univérso mundo.}
	\\
	\instruct{Hymnus -- Crux fidelis}
	\markup{Ant.} Crux fidélis, inter omnes arbor una nóbilis: nulla silva talem profert fronde, flore, gérmine. * Dulce lignum dulces clavos, dulce pondus sústinet.
	\\\\
	Pange, lingua, gloriósi
	láuream certáminis,
	et super Crucis trophæo
	dic triúmphum nóbilem:
	quáliter Redémptor orbis
	immolátus vícerit
	\\\\
	\markup{Ant.} Crux fidélis, inter omnes arbor una nóbilis: nulla silva talem profert fronde, flore, gérmine.
	\\\\
	De paréntis protoplásti
	fraude Factor cóndolens,
	quando pomi noxiális
	in necem morsu ruit:
	ipse lignum tunc notávit,
	damna ligni ut sólveret.
	\\\\
	\markup{Ant.} Dulce lignum dulces clavos, dulce pondus sústinet.
	\\\\
	Hoc opus nostræ salútis
	ordo depopóscerat:
	multifórmis proditóris
	ars ut artem fálleret:
	et medélam ferret inde,
	hostis unde læserat.
	\\\\
	\markup{Ant.} Crux fidélis, inter omnes arbor una nóbilis: nulla silva talem profert fronde, flore, gérmine.
	\\\\
	Quando venit ergo sacri
	plenitúdo témporis,
	missus est ab arce Patris
	Natus, orbis Cónditor:
	atque ventre virgináli
	carne amíctus pródiit.
	\\\\
	\markup{Ant.} Dulce lignum dulces clavos, dulce pondus sústinet.
	\\\\
	Vagit Infans inter arcta
	cónditus præsépia:
	membra pannis involúta
	Virgo Mater álligat:
	et Dei manus pedésque
	stricta cingit fáscia.
	\\\\
	\markup{Ant.} Crux fidélis, inter omnes arbor una nóbilis: nulla silva talem profert fronde, flore, gérmine.
	\\\\
	Lustra sex qui jam perégit,
	tempus implens córporis,
	sponte líbera Redémptor
	passióni déditus,
	Agnus in Crucis levátur
	immolándus stípite.
	\\\\
	\markup{Ant.} Dulce lignum dulces clavos, dulce pondus sústinet.
	\\\\
	Felle potus ecce languet:
	spina, clavi, láncea
	mite corpus perforárunt,
	unda manat et cruor:
	terra, pontus, astra, mundus,
	quo lavántur flúmine!
	\\\\
	\markup{Ant.} Crux fidélis, inter omnes arbor una nóbilis: nulla silva talem profert fronde, flore, gérmine.
	\\\\
	Flecte ramos, arbor alta,
	tensa laxa víscera,
	et rigor lentéscat ille,
	quem dedit natívitas:
	et supérni membra Regis
	tende miti stípite.
	\\\\
	\markup{Ant.} Dulce lignum dulces clavos, dulce pondus sústinet.
	\\\\
	Sola digna tu fuísti
	ferre mundi víctimam:
	atque portum præparáre
	arca mundo náufrago:
	quam sacer cruor perúnxit,
	fusus Agni córpore.
	\\\\
	\markup{Ant.} Crux fidélis, inter omnes arbor una nóbilis: nulla silva talem profert fronde, flore, gérmine.
	\\\\
	Sempitérna sit beátæ
	Trinitáti glória:
	æqua Patri Filióque;
	par decus Paráclito:
	Uníus Triníque nomen
	laudet univérsitas.
	Amen.
	\\\\
	\markup{Ant.} Dulce lignum dulces clavos, dulce pondus sústinet.
	\\\\
	\instruct{Circa finem adorationis Crucis accenduntur candelæ super Altare: et Diaconus accipiens Bursam Corporalium, extendit Corporale more solito, et juxta illud ponit Purificatorium: et finita adoratione, Crucem reverenter accipit, et reportat ad Altare. Postea ordinatur Processio ad locum, ubi pridie Sacramentum repositum fuerat.
	\\\\
	Præcedit Subdiaconus cum Cruce inter duos Acolythos deferéntes candelábra cum cereis accensis, et Clerus per ordinem, ultimus Sacerdos cum Ministris. Cum perventum fuerit ad locum Sacramenti, accenduntur intorticia, quæ non exstinguuntur nisi post sumptionem Sacramenti: et Sacerdos genuflectit ante Sacramentum, orat aliquantulum: Diaconus interim aperit capsulam, in qua reconditum est Corpus Domini: postea Sacerdos surgens, absque benedictione imponit incensum in duobus thuribulis, Diacono naviculam ministrante, et genuflexus incensat Sacramentum: tum Diaconus extrahens Calicem cum Sacramento de capsula, dat ad manus Sacerdotis, et tegit extremitatibus Veli, quod húmeros illius circumdat: et procedunt ordine quo venerunt: defertur baldachinum super Sacramentum: et duo Acolythi cum thuribulis continue Sacramentum incensant: interim cantatur Hymnus Vexilla Regis pródeunt.}
	\\
	\instruct{Hymnus -- Vexilla Regis pródeunt}
	Vexílla Regis pródeunt;
	Fulget crucis mystérium,
	Qua vita mortem pértulit,
	Et morte vitam prótulit.
	\\\\
	Quæ vulneráta lánceæ
	Mucróne diro, críminum
	Ut nos laváret sórdibus,
	Manávit unda et sánguine.
	\\\\
	Impléta sunt quæ cóncinit
	David fidéli cármine,
	Dicéndo natiónibus:
	Regnávit a ligno Deus.
	\\\\
	Arbor decóra et fúlgida,
	Ornáta regis púrpura,
	Elécta digno stípite
	Tam sancta membra tángere.
	\\\\
	Beáta, cujus bráchiis
	Prétium pepéndit sæculi,
	Statéra facta córporis,
	Tulítque prædam tártari.
	\\\\
	O Crux, ave, spes única,
	Paschále quæ fers gáudium,
	Piis adáuge grátiam,
	Reísque dele crímina.
	\\\\
	Te, fons salútis, Trínitas,
	Colláudet omnis spíritus:
	Quibus crucis victóriam
	Largíris, adde præmium.
	Amen.
	
	\section{Communio}
	
	\instruct{Cum venerit Sacerdos ad Altare, posito super illud Calice, genuflexus rursum incensat, et accedens deponit Hostiam ex Calice super Patenam, quam Diaconus tenet: et accipiens Patenam de manu Diaconi, Hostiam sacram ponit super Corporale, nihil dicens. Si tetigerit Sacramentum, digitos abluat in aliquo vase. Interim Diaconus imponit vinum in Calicem, et Subdiaconus aquam, quam Sacerdos non benedicit, nec dicit super eam orationem consuetam: sed accipiens Calicem a Diacono, ponit super Altare, nihil dicens: et Diaconus illum cooperit Palla: deinde imponit incensum in thuribulo absque benedictione, et incensat Oblata, Crucem et Altare more solito, genuflectens ante et post, et quandocumque transit ante Sacramentum. Cum incensat Oblata, dicit:}
	\\
	\priest{Incénsum istud, a te benedíctum, ascéndat ad te, Dómine: et descéndat super nos misericórdia tua.}
	\\
	\instruct{Cum incensat Altare, dicit:}
	\\
	\instruct{Ps 140:2-4.}
	\priest{Dirigátur, Dómine, orátio mea, sicut incénsum in conspéctu tuo: elevátio mánuum meárum sacrifícium vespertínum. Pone, Dómine, custódiam ori meo, et óstium circumstántiæ lábiis meis: ut non declínet cor meum in verba malítiæ, ad excusándas excusatiónes in peccátis.}
	\\
	\instruct{Quando reddit thuribulum Diacono, dicit:}
	\\
	\priest{Accéndat in nobis Dóminus ignem sui amóris, et flammam ætérnæ cantátis. Amen.}
	\\
	\instruct{Et ipse non incensatur.}
	\\
	\instruct{Postea aliquantulum extra Altare in cornu Epistolæ lavat manus, nihil dicens: deinde in medio Altaris inclinatur, junctis manibus, dicit:}
	\\
	\priest{In spíritu humilitátis et in ánimo contríto suscipiámur a te, Dómine: et sic fiat sacrifícium nostrum in conspéctu tuo hódie, ut pláceat tibi, Dómine Deus.}
	\\
	\instruct{Deinde versus ad populum in cornu Evangelii, dicit more solito:}
	\\
	\priest{Oráte, fratres, ut meum ac vestrum sacrifícium acceptábile fiat apud Deum Patrem omnipoténtem.}
	\\
	\instruct{Et per eamdem viam revertitur, non perficiens circulum: et consequenter, omissis aliis, dicit:}
	\\
	\priest{Orémus: Præcéptis salutáribus móniti, et divína institutióne formáti, audémus dícere:
	\\\\
	Pater noster, qui es in cœlis: Sanctificétur nomen tuum: Advéniat regnum tuum: Fiat volúntas tua, sicut in cœlo, et in terra. Panem nostrum quotidiánum da nobis hódie: Et dimítte nobis débita nostra, sicut et nos dimíttimus debitóribus nostris. Et ne nos indúcas in tentatiónem.}
	\server{Sed líbera nos a malo.}
	\\
	\instruct{Sacerdos, sub silentio dicto:}
	\\
	\priest{Amen.}
	\\
	\instruct{Eadem voce, qua dixit Pater noster, absolute sine Orémus in tono orationis Missæ ferialis dicit:}
	\\
	\priest{Líbera nos, quæsumus, Dómine, ab ómnibus malis, prætéritis, præséntibus et futúris: et intercedénte beáta et gloriósa semper Vírgine Dei Genitríce María, cum beátis Apóstolis tuis Petro et Paulo, atque Andréa, et ómnibus Sanctis, da propítius pacem in diébus nostris: ut, ope misericórdiæ tuæ adjúti, et a peccáto simus semper líberi et ab omni perturbatióne secúri. Per eúndem Dóminum nostrum Jesum Christum, Fílium tuum: Qui tecum vivit et regnat in unitáte Spíritus Sancti Deus, per ómnia sæcula sæculórum.}
	\priest{Amen.}
	\\
	\instruct{Tunc Celebrans, facta reverentia usque ad terram, supponit Patenam Sacramento, quod in dextera accipiens, elevat, ut videri possit a populo: et statim supra Calicem dividit in tres partes, quarum ultimam mittit in Calicem more solito, nihil dicens. Pax Dómini non dicitur, nec Agnus Dei, neque pacis osculum datur. Postmodum, prætermissis duabus primis orationibus, dicit tantum sequentem:}
	\\
	Percéptio Córporis tui, Dómine Jesu Christe, quod ego indígnus súmere præsúmo, non mihi provéniat in judícium et condemnatiónem: sed pro tua pietáte prosit mihi ad tutaméntum mentis et córporis, et ad medélam percipiéndam: Qui vivis et regnas cum Deo Patre in unitáte Spíritus Sancti Deus, per ómnia sæcula sæculórum. Amen.
	\\\\
	\instruct{Tunc genuflectit, et accipit Patenam cum Corpore Christi: et maxima humilitate ac reverentia dicit:}
	\\
	Panem cœléstem accípiam, et nomen Dómini invocábo.
	\\\\
	\instruct{Percutit pectus suum, ter dicens:}
	\\
	Dómine, non sum dignus, ut intres sub tectum meum: sed tantum dic verbo, et sanábitur ánima mea.
	\\
	Dómine, non sum dignus, ut intres sub tectum meum: sed tantum dic verbo, et sanábitur ánima mea.
	\\
	Dómine, non sum dignus, ut intres sub tectum meum: sed tantum dic verbo, et sanábitur ánima mea.
	\\\\
	\instruct{Postea signat se Sacramento, dicens:}
	\\
	Corpus Dómini nostri Jesu Christi custódiat ánimam meam in vitam ætérnam. Amen.
	\\\\
	\instruct{Et sumit Corpus reverenter.}
	\\
	\instruct{Deinde, omissis omnibus, quæ dici solent ante sumptionem Sanguinis, immediate particulam Hostiæ cum vino reverenter sumit de Calice. Et more solito facta ablutione digitorum, et sumpta purificatione, in medio Altaris inclinatus, manibus junctis, dicit:}
	\\
	Quod ore súmpsimus, Dómine, pura mente capiámus: et de múnere temporáli fiat nobis remédium sempitérnum.
	\\\\
	\instruct{Non dicitur Corpus tuum, Dómine, nec Postcommunio, nec Pláceat tibi, nec datur benedictio; sed facta reverentia Altari, Sacerdos cum Ministris discedit: et dicuntur Vesperæ sine cantu: et denudatur Altare.}
	
	\newpage 
	
	\begin{flushleft}
		\textbf{Missa Præsanctificatorum} -- \textit{Mass of the Presanctified}
	\end{flushleft}
	
	\begin{flushleft}
		Copyright © Chavez Cyrillus Mariae Poon 2026. All rights reserved.
	\end{flushleft}
	
	\begin{flushleft}
		This document is indented for internal distribution only.
	\end{flushleft}
	
	\begin{flushleft}
		It may be freely shared, provided that the source is properly attributed, and it is not used for commercial purposes.
	\end{flushleft}
	
	\begin{flushleft}
		No modification or resale of the content are permitted without prior written permission.
	\end{flushleft}
	
	\begin{flushleft}
		This work is free of known copyright restrictions.
	\end{flushleft}
	
	\begin{flushleft}
		For more information, please visit: https://www.chavezpoon.com
	\end{flushleft}
	
	\begin{flushleft}
		Typesetting using \LaTeX and liturg package.
	\end{flushleft}
	\newpage
\end{document}
